%% 第 13 章：Jordan 块与 Jordan 标准型
%% 本文件由 tools/md2tex.py 自动生成，请勿直接编辑。

\chapter{Jordan 块与 Jordan 标准型}


\begin{jdefn}{Jordan 块}{r:14:1}
对长度为 $r$ 的 Jordan 链：

\[
(v_1,\ldots,v_r),
\]

其中：

\[
Tv_1=\lambda v_1,
\]

\[
Tv_j=\lambda v_j+v_{j-1},
\]

则 $T$ 在这组基下的矩阵为：

\[
J_r(\lambda)
=
\begin{pmatrix}
\lambda&1&0&\cdots&0\\
0&\lambda&1&\cdots&0\\
0&0&\lambda&\ddots&0\\
\vdots&\vdots&\ddots&\ddots&1\\
0&0&0&0&\lambda
\end{pmatrix}.
\]

它称为一个属于特征值 $\lambda$ 的 Jordan 块。
\end{jdefn}

\section{Jordan 块的意义}

对角线上的 $\lambda$ 表示该链上的基本缩放。

上方的 $1$ 表示：这个方向除了被缩放外，还会向链中更低一级的方向产生不可消除的耦合。

若块大小大于 $1$，该部分不能仅由普通特征向量组成，因此不能完全对角化。

\begin{jthm}{Jordan 标准型定理}{r:14:2}
若 $V$ 是有限维复向量空间，则任意算子：

\[
T:V\to V
\]

都存在一组基，使其矩阵是若干 Jordan 块构成的块对角矩阵：

\[
[T]
=
J_{r_1}(\lambda_1)
\oplus
\cdots
\oplus
J_{r_m}(\lambda_m).
\]

Jordan 块的排列顺序可以改变，但每个特征值对应的块大小集合是由 $T$ 唯一确定的。
\end{jthm}

\begin{jprop}{Jordan 标准型与对角化}{r:14:3}
\[
T\text{ 可对角化}
\iff
\text{所有 Jordan 块都为 }1\times1.
\]

等价地：

\[
G_\lambda=E_\lambda
\]

对每个特征值 $\lambda$ 都成立。

还等价于最小多项式没有重复根：

\[
m_T(z)
=
\prod_\lambda(z-\lambda).
\]
\end{jprop}

\section{幂零算子的矩阵结构}

若 $N$ 是幂零算子，则它唯一的特征值是：

\[
0.
\]

证明：设 $\lambda$ 是 $N$ 的任意一个特征值，$v\neq 0$ 为对应的特征向量，则

\[
Nv=\lambda v.
\]

两边连续作用算子 $N$，可得

\[
N^2 v = N(Nv)=N(\lambda v)=\lambda Nv=\lambda^2 v.
\]

归纳可得

\[
N^k v=\lambda^k v.
\]

由幂零条件 $N^k=0$，于是

\[
\lambda^k v = N^k v = 0.
\]

因 $v\neq 0$，故 $\lambda^k=0$，从而 $\lambda=0$。

因此幂零算子仅能以 $0$ 作为特征值。

因此，它在 Jordan 基下的矩阵由若干：

\[
J_r(0)
\]

构成。

这些 Jordan 块的对角线全为零。

显然，对于使其成为上三角阵的任意一组基，对角线为零。

幂零算子的矩阵并非在任意基下都显然对角线为零；这里的结论是针对 Jordan 基，或任意使其上三角化的基。
